\documentclass[11pt]{article}
\usepackage[margin=1in]{geometry}
\usepackage{amsmath,amssymb,amsthm}
\usepackage{graphicx}
\usepackage{hyperref}

\newtheorem{theorem}{Theorem}
\newtheorem{corollary}{Corollary}
\newtheorem{lemma}{Lemma}

\DeclareMathOperator{\Ext}{Ext}
\DeclareMathOperator{\supp}{supp}

\title{Extreme Points of Positive-Definite Extension Sets}
\author{Agent lane 1}
\date{July 18, 2026}

\begin{document}
\maketitle

\section*{Status}

Candidate full solution, likely valid, for the locally compact Abelian measure-valued formulation of Question 3.1 in Jorgensen--Pedersen--Tian \cite{JPT}. The result is an abstract characterization rather than a closed-form parametrization.

\section*{Source Problem}

The source paper studies positive definite extensions of a continuous positive definite function \(F\) initially defined on a difference set. On page 65 it asks:

\begin{center}
\includegraphics[width=0.94\linewidth]{figures/open_problem_crop.png}
\end{center}

In the locally compact Abelian case, using the source's Bochner-measure model, let \(G\) be the ambient group, let \(\widehat G\) be its dual, and put
\[
  D=\Omega^{-1}\Omega .
\]
After the source's normalization \(F(e)=1\), write
\[
  \Ext(F)
  =
  \left\{\mu\in \mathcal P(\widehat G):
  \widehat\mu(x)=F(x)\ \text{for every }x\in D
  \right\},
\]
where
\[
  \widehat\mu(x)=\int_{\widehat G} \langle \lambda,x\rangle\,d\mu(\lambda).
\]

\section*{Proof Intuition}

Convex decompositions of a measure \(\mu\) inside \(\Ext(F)\) are exactly bounded signed perturbations of \(\mu\) whose Fourier transform vanishes on the observed set \(D=\Omega^{-1}\Omega\). Indeed, if \(\mu=t\mu_1+(1-t)\mu_2\), then \(\mu_1\) and \(\mu_2\) are automatically absolutely continuous with respect to \(\mu\), with bounded Radon--Nikodym derivatives. Conversely, any bounded real density \(h\) invisible to all observed characters produces the two distinct extension measures \((1\pm\varepsilon h)\mu\). This gives the full characterization. In the Euclidean compact-support subcase, invisible perturbations vanish by analytic continuation, giving a concrete sufficient condition for extremality.

\section*{The Extremality Characterization}

\begin{theorem}[bounded perturbation characterization]
\label{lem:criterion}
Let \(G\) be a locally compact Abelian group, let \(\Omega\subset G\) be nonempty, let \(D=\Omega^{-1}\Omega\), and let \(F:D\to\mathbb C\) be continuous positive definite with \(F(e)=1\). Assume \(\Ext(F)\neq\emptyset\), where \(\Ext(F)\) is identified with the probability measures \(\mu\) on \(\widehat G\) satisfying \(\widehat\mu|_D=F\). For \(\mu\in\Ext(F)\), the following are equivalent:
\begin{enumerate}
\item \(\mu\) is an extreme point of \(\Ext(F)\).
\item There is no nonzero real-valued \(h\in L^\infty(\mu)\) such that
\[
  \int_{\widehat G} \langle \lambda,x\rangle h(\lambda)\,d\mu(\lambda)=0,
  \qquad x\in D.
\]
\item The real linear span of
\[
  \left\{\operatorname{Re}\langle \cdot,x\rangle,\operatorname{Im}\langle \cdot,x\rangle : x\in D\right\}
\]
is dense in \(L^1(\mu;\mathbb R)\).
\end{enumerate}
\end{theorem}

\begin{proof}
We prove first the equivalence of (1) and (2). Suppose that \(\mu\) is not extreme. Then
\[
  \mu=t\mu_1+(1-t)\mu_2
\]
for some \(0<t<1\) and distinct \(\mu_1,\mu_2\in\Ext(F)\). Since \(\mu_i\ll\mu\), write \(g_i=d\mu_i/d\mu\). The identity above gives \(tg_1+(1-t)g_2=1\), and in particular \(g_i\in L^\infty(\mu)\). Since both \(\mu_i\) and \(\mu\) have the same Fourier transform on \(D\),
\[
  \int_{\widehat G}\langle \lambda,x\rangle(g_1(\lambda)-1)\,d\mu(\lambda)=0,
  \qquad x\in D.
\]
The real function \(h=g_1-1\) is not zero, otherwise \(\mu_1=\mu\) and then also \(\mu_2=\mu\).

Conversely, let \(h\in L^\infty(\mu)\) be real, nonzero, and annihilate all characters indexed by \(D\). Since \(0\in D\), taking \(x=0\) gives \(\int h\,d\mu=0\). Choose \(\varepsilon>0\) so small that \(1\pm \varepsilon h\ge 0\) \(\mu\)-almost everywhere, and define
\[
  d\mu_\pm=(1\pm \varepsilon h)\,d\mu.
\]
Then \(\mu_\pm\) are probability measures, \(\mu=(\mu_++\mu_-)/2\), and their Fourier transforms agree with \(\widehat\mu=F\) on \(D\). Hence \(\mu_\pm\in\Ext(F)\), and the decomposition is nontrivial.

Finally, (2) and (3) are the standard annihilator reformulation. The annihilator in \(L^\infty(\mu;\mathbb R)\) of the real span of the observed characters is precisely the set of real \(h\in L^\infty(\mu)\) satisfying the displayed equations in (2). Thus (2) says that this annihilator is zero. By the Hahn--Banach separation theorem, a real linear subspace of \(L^1(\mu;\mathbb R)\) has zero annihilator in \(L^\infty(\mu;\mathbb R)\) if and only if it is norm dense.
\end{proof}

\section*{Euclidean Compact-Support Consequences}

\begin{corollary}
\label{thm:compact-extreme}
Let \(\Omega\subset\mathbb R^n\) be nonempty and open, and let \(F:\Omega-\Omega\to\mathbb C\) be continuous and positive definite. If \(\mu\in\Ext(F)\) has compact support in \(\mathbb R^n\), then \(\mu\) is an extreme point of \(\Ext(F)\).
\end{corollary}

\begin{proof}
By Lemma \ref{lem:criterion}, it suffices to show that no nonzero bounded perturbation can have a Fourier transform vanishing on \(D=\Omega-\Omega\). Let \(h\in L^\infty(\mu)\) be real-valued and suppose
\[
  \int e^{i x\cdot\lambda}h(\lambda)\,d\mu(\lambda)=0,
  \qquad x\in D.
\]
Set \(d\sigma=h\,d\mu\). Since \(\mu\) has compact support, \(\sigma\) is a compactly supported finite signed measure. Its Fourier transform extends to an entire function on \(\mathbb C^n\):
\[
  \Phi(z)=\int_{\mathbb R^n} e^{i z\cdot\lambda}\,d\sigma(\lambda),
  \qquad z\in\mathbb C^n.
\]
The set \(D\) is a nonempty open subset of \(\mathbb R^n\), and \(\Phi\) vanishes on \(D\). Therefore the real-analytic function \(\Phi|_{\mathbb R^n}\) has all derivatives zero at every point of \(D\), so the entire power series of \(\Phi\) at any such point is zero. Hence \(\Phi\equiv 0\) on \(\mathbb C^n\).

In particular, the Fourier transform of \(\sigma\) vanishes on all of \(\mathbb R^n\). By uniqueness of the Fourier transform for finite Borel measures on \(\mathbb R^n\), \(\sigma=0\). Thus \(h=0\) in \(L^\infty(\mu)\), and Lemma \ref{lem:criterion} implies that \(\mu\) is extreme.
\end{proof}

\begin{corollary}
For a fixed \(F\) and \(\Omega\subset\mathbb R^n\), the set \(\Ext(F)\) contains at most one compactly supported measure.
\end{corollary}

\begin{proof}
If \(\mu,\nu\in\Ext(F)\) are compactly supported, then \(\sigma=\mu-\nu\) is compactly supported and \(\widehat\sigma=0\) on \(\Omega-\Omega\). The same analytic-continuation and Fourier-uniqueness argument gives \(\sigma=0\).
\end{proof}

\section*{Limitations And Novelty Check}

The main theorem is a full abstract characterization of the extreme boundary of the measure-valued \(\Ext(F)\) in the locally compact Abelian setting of Question 3.1. It is not a parametrization by a small named family of measures, and it does not classify non-Abelian unitary-representation triples from the source's broader \(Ext(F)\) formalism. The compact-support corollary uses Euclidean analytic continuation and is not asserted for arbitrary locally compact Abelian groups.

The cheap run indexes were searched for arXiv:1507.02547, the source title, positive-definite extension keywords, and extreme-point keywords. No exact prior packet was found before this work. Web searches for exact and close phrases around \(\Ext(F)\), positive definite extensions, bounded perturbations, \(L^\infty\)-criteria, Radon--Nikodym derivatives, compact support, and extreme points found the source paper and neighboring positive-definite-function literature, but no exact statement of the characterization above.

\section*{Human Review Recommendation}

Review as a full abstract answer to Question 3.1 in the locally compact Abelian measure formulation. The main check is that this is the intended reading of the source question; under that reading, the theorem gives necessary and sufficient conditions for membership in \(ext(\Ext(F))\).

\begin{thebibliography}{9}

\bibitem{JPT}
Palle Jorgensen, Steen Pedersen, and Feng Tian,
\emph{Extensions of Positive Definite Functions: Applications and Their Harmonic Analysis},
arXiv:1507.02547, 2015.

\bibitem{Rudin}
Walter Rudin,
\emph{Fourier Analysis on Groups},
Interscience Publishers, 1962.

\end{thebibliography}

\end{document}
